\section{An obstruction to every equivariant smooth cutoff transport}
\label{sec:equivariant-cutoff-obstruction}

\paragraph{Status.}
The complete post-writing author reading and two independent
whole-proof readings have passed.  The separate source-interface
readings retain their expressly stated scope.
The hypotheses specify an actual class of maps.  No assertion is
made about arbitrary non-equivariant maps or a changed cutoff family.

The preceding MA-PI-001--004 refute the canonical theta map.
Here the map is allowed to vary.  The obstruction concerns exact
orthogonal-cutoff transport at all sufficiently large levels,
already on the two constant plateaus of the physical cutoff.
The smooth-source assertion uses the full classical zero count
and containing critical strip, without placing zeros in a
particular narrow strip.  A separate,
weaker-regularity assertion explicitly uses known critical-line zeros.

Put $\mathcal D=C_c^\infty(\mathbb R)$ and
$\tau_yf(x)=f(x-y)$.  Let $\mathcal E$ be the original
all-exponential Schwartz space, with seminorms
$\sup_x e^{N|x|}|f^{(j)}(x)|$ for every $N,j$.
Its Mellin transform is
$\mathcal Mf(z)=\int_{\mathbb R}f(x)e^{zx}\,dx$.
Fix the actual symmetric Hilbert metric
\[
 H_{a,s}=L^2(\mathbb R,e^{2a|x|}(1+x^2)^{s/2}\,dx),
 \qquad a>0,\quad s\ge0,
 \qquad b=-a.
 \tag{MA-PI.150}
\]
The topology on $\mathcal D$ is its usual strict inductive
limit topology.  Continuity into $H_{a,s}$, or into
$\mathcal E$, always refers to the indicated actual topology.

\paragraph{MA-PI-005 (no nonzero equivariant smooth transport with finite Gram propagation).}
Let $\mathscr I\subset\mathcal E$ be the original rational
smooth theta ideal.  It is enough here to use its proved
full-jet annihilation property
\[
 (\mathcal Mh)^{(j)}(\rho-1/2)=0
 \quad(h\in\mathscr I,\ 0\le j<m_\rho),
 \tag{MA-PI.151}
\]
for every nontrivial zeta zero $\rho$.
Suppose a continuous linear map $T:\mathcal D\to\mathcal E$
has image in $\mathscr I$ and commutes with every $\tau_y$.
Assume there is a finite $R\ge0$ such that, whenever the compact
supports of $f,g\in\mathcal D$ are ordered and separated by
more than $R$, one has
\[
 \langle T\tau_{-L}f,T\tau_{-L}g\rangle_{H_{a,s}}=0
 \quad\hbox{for all sufficiently large $L$}.
 \tag{MA-PI.152}
\]
The threshold for $L$ may depend on $f,g$.  Then $T=0$.
The same conclusion holds for a continuous equivariant map
defined on
\[
 \mathcal D_0=\left\{v\in\mathcal D:
                  \int e^{-x/2}v(x)\,dx=0\right\},
 \tag{MA-PI.153}
\]
if (MA-PI.152) holds for its vectors.

\begin{proof}
\emph{1. Extracting the actual one-sided Gram form.}
For $h,k\in H_{a,s}$, a change of variables and dominated
convergence give
\[
 \lim_{L\to\infty}e^{-2aL}L^{-s}
       \langle\tau_{-L}h,\tau_{-L}k\rangle_{H_{a,s}}
       =\int h(x)\overline{k(x)}e^{-2ax}\,dx.
 \tag{MA-PI.154}
\]
Indeed the translated weight divided by $e^{2aL}L^s$
tends to $e^{-2ax}$ and, for $L\ge1$, is bounded by
$C_s e^{2a|x|}(1+x^2)^{s/2}$.  Cauchy--Schwarz supplies
an integrable majorant for the cross term.  Thus (MA-PI.152)
implies orthogonality in $L^2(e^{2bx}dx)$ for every such pair.

Define, on $\mathcal D$,
\[
 Af(x)=e^{bx}T(e^{-b\,\cdot}f)(x).
 \tag{MA-PI.155}
\]
This is continuous into ordinary $L^2$, commutes with
translations, and preserves the separated-support
orthogonality just obtained.  It extends continuously
$\mathcal S\to L^2$: take a fixed smooth partition on
unit intervals, translate each piece back to one fixed
compact interval, and use the one finite $C^N$ seminorm
bound there.  The sum of those bounds is controlled by
a Schwartz seminorm with two additional powers of decay.
The series of outputs converges absolutely in $L^2$,
agrees with $A$ on $\mathcal D$, and preserves translation
equivariance by density.

\emph{2. A compact Gram distribution.}
Use the Fourier convention
$\widehat f(t)=\int f(x)e^{itx}\,dx$, with the inverse
factor $(2\pi)^{-1}$.
For $f,g\in\mathcal S$, equivariance and the convergent
convolution integral give
$A(f*g)=(Af)*g=f*(Ag)$.
Choose a Gaussian $\varphi$ whose Fourier transform never
vanishes and define the locally square-integrable function
$M=\widehat{A\varphi}/\widehat\varphi$.
Taking the Fourier transform of the displayed convolution
identity proves
\[
 \widehat{Af}(t)=M(t)\widehat f(t),\qquad
 \langle Af,Ag\rangle
       =\frac{1}{2\pi}\int |M(t)|^2
                   \widehat f(t)\overline{\widehat g(t)}\,dt.
 \tag{MA-PI.156}
\]
Continuity on Schwartz space, applied to the modulated
Gaussian, bounds the $L^2$ mass of $M$ on every unit
frequency interval by one fixed polynomial in its center.
Consequently $|M|^2$ is a tempered distribution.
Let $k$ be the tempered distribution whose Fourier
transform is $|M|^2$, with this specified convention.
The separated-support orthogonality implies
\[
                       \operatorname{supp}k\subset[-R,R].
 \tag{MA-PI.157}
\]
To verify that support assertion directly, put one test
factor in a shrinking neighborhood of zero and the other
in either component of $\mathbb R\setminus[-R,R]$.
Their correlation tends to the latter test function (with
the harmless reversal prescribed by the Fourier convention).
Formula (MA-PI.156) is zero for every sufficiently small
neighborhood.  Taking the distributional limit proves
that $k$ vanishes on both components.

The Fourier transform of this compact distribution has
an entire extension $E(w)$, agreeing with $|M(t)|^2$ on
the real line as distributions, and obeying, for some
$C,N$ and finite $R_1$,
\[
 |E(w)|\le C(1+|w|)^N e^{R_1|\Im w|}.
 \tag{MA-PI.158}
\]
For example evaluate $k$ on a compact cutoff times
$e^{iwx}$; its finite distributional order gives this
bound directly.  The cutoff may have support in a fixed
slightly larger interval than $[-R,R]$.  If $E\ne0$,
the Jensen estimate already used in (MA-MO.412) gives
only $O(T)$ zeros in $|w|\le T$, counted with multiplicity.

\emph{3. The genuine entire Mellin multiplier.}
The original $T$, not its Schwartz extension in Step 1,
has an entire Mellin multiplier $m(z)$.
In fact, for $f,g\in\mathcal D$,
$T(f*g)=(Tf)*g=(Tg)*f$ in $\mathcal E$, so
\[
     \mathcal M(Tf)(z)\mathcal Mg(z)
           =\mathcal M(Tg)(z)\mathcal Mf(z).
 \tag{MA-PI.159}
\]
At each $z$ choose $g$ with $\mathcal Mg(z)\ne0$.
The local quotients $\mathcal M(Tg)/\mathcal Mg$ patch,
by (MA-PI.159), to an entire function satisfying
$\mathcal M(Tf)=m\mathcal Mf$ for every $f\in\mathcal D$.
There is no meromorphic pole at a zero of one chosen
denominator: another compact test gives a nonzero
denominator there.  Formula (MA-PI.155) identifies
$M(t)=m(b+it)$ almost everywhere.  Since the latter is
continuous, Step 2 yields the pointwise identity
\[
 m(z)\overline{m(2b-\bar z)}
                  =E\bigl(-i(z-b)\bigr)\qquad(z\in\mathbb C).
 \tag{MA-PI.160}
\]
Both sides are entire and agree on the vertical line
$\Re z=b$, which proves the identity everywhere.

For every nontrivial zeta zero $\rho$, choose a compact
test whose Mellin value at $\rho-1/2$ is nonzero.
The chosen input is nonzero in Mellin coordinates at that
point, so the full-jet condition (MA-PI.151) forces
$\operatorname{ord}_{\rho-1/2}m\ge m_\rho$.
Thus $E$ has all the corresponding
zeros at
$-i(\rho-1/2-b)$.
The classical unconditional count, in the exact form
Greni\'e--Molteni (2019), Appendix A, Theorem A.1,
author p.~26, is of order $T\log T$ for $|\Im\rho|\le T$,
with multiplicity.  Specialize to the trivial character
of $\mathbb Q$, as in GRH\_NIGHT NT-R08c.  The zeros lie
in the classical critical strip, so these transformed
points have modulus at most $T+O_a(1)$.
All those multiplicities are retained in the count.
No assertion about their position inside a narrower strip
has been used.
This contradicts the $O(T)$ bound unless $E=0$.
Then (MA-PI.156) implies $A=0$, hence $T=0$.

\emph{4. The rational polar moment.}
The differential map $D-1/2$ sends $\mathcal D$ onto
$\mathcal D_0$, preserves supports, and commutes with
translations.  Surjectivity is explicit: for
$v\in\mathcal D_0$ put
\[
 u(x)=e^{x/2}\int_{-\infty}^x e^{-t/2}v(t)\,dt.
 \tag{MA-PI.161}
\]
The vanishing moment makes $u$ compactly supported and
$(D-1/2)u=v$.  Apply the preceding proof to
$T(D-1/2)$; its image and separated-support properties
are unchanged.  Its vanishing and surjectivity give $T=0$.
\end{proof}

\paragraph{Actual rational cutoff coordinates and the scope of the corollary.}
For an even physical Fourier input $u$, smooth and compactly
supported away from zero on the positive half-line, put
\[
 v(x)=e^{-x/2}u(e^{-x}),\qquad f=\mathcal F_\infty^{-1}u.
 \tag{MA-PI.162}
\]
The even extension of $u$ is understood.  These are genuine
even Schwartz inputs.  Their two original polar conditions
are $\int f=u(0)=0$ and
$f(0)=2\int_0^\infty u(\xi)d\xi=0$; the second is
exactly $v\in\mathcal D_0$.
These are genuine inputs of the original rational global
PW source: Masullo, Definition~7.3.2.3 and (7.3.2.12)--(7.3.2.17),
pp.~354--355, with $S=\{\infty\}$.  They are not thereby
identified with the coordinate-flat local source on p.~515:
$f(0)=0$ alone does not make $f$ flat at zero.
If $v$ is replaced by $\tau_yv$, then
$u_y(\xi)=e^{y/2}u(e^y\xi)$ and
$f_y(t)=e^{-y/2}f(e^{-y}t)$.  Thus the actual centered
theta normalization (MA-PI.2), with the principal
character and its two moments, changes its global
logarithmic image by precisely $\tau_y$.

Let $P_Y$ be the actual orthogonal global PW projection
in the fixed metric (MA-PI.150), as in MA-GP and MA-PI-003.
Any continuous equivariant smooth-source replacement
$T:\mathcal D_0\to\mathscr I$ that intertwines the
physical cutoffs with these projections must, at a minimum,
satisfy their plateau relations
\[
 \begin{array}{ll}
 P_YTv=Tv,&\operatorname{supp}v\subset[c+\log2,\infty),\\
 P_YTv=0,&\operatorname{supp}v\subset(-\infty,c-\log2],
 \end{array}
 \qquad c=-\log Y.
 \tag{MA-PI.163}
\]
These are the ball-cutoff plateau assumptions used in the
physical Fourier construction of Masullo, p.~516, Step~II,
and its p.~519 trace application (not the separate annular
cutoff of Definition~9.3.3.1):
$\rho_Y(r)=1$ for $r\le Y/2$ and $\rho_Y(r)=0$ for
$r\ge2Y$, since the physical radial Fourier coordinate
is $r=e^{-x}$.  They hold eventually for the source's
thin transition around the physical frequency ball.
Only plateau inputs are used, so no
claim is needed that multiplication by an arbitrary
transition profile preserves $\mathcal D_0$.

For supports separated by more than $2\log2$, choose
$c$ between them.  The same orthogonal $P_Y$ fixes
one output and kills the other, so the outputs are
orthogonal.  Translating both inputs far to the left
makes $Y$ arbitrarily large.  Thus (MA-PI.163), even
only for all sufficiently large $Y$, implies
(MA-PI.152) with $R=2\log2$.
MA-PI-005 proves that every such $T$ is zero.
In particular, adding an arbitrary continuous
equivariant smooth correction to the theta map
cannot yield a nonzero map with these exact plateau
and orthogonal-projection identities.

This conclusion is stronger than a counterexample to
one prescribed map.  Its precise range is still
important: it concerns smooth ideal-valued maps,
equivariance, and this fixed Hilbert metric and cutoff
family.  It does not assert that the draft printed a
particular intertwining equation, that a proposed
global-to-local map necessarily has such a smooth
inverse, or that other trace comparisons are impossible.
The source-interface question is distinct from the
proved obstruction to the explicitly specified maps.

\paragraph{MA-PI-006 (the Hilbert-target variant and its additional input).}
Let $J\subset H_{a,s}$ be the actual closed theta ideal.
Suppose $T:\mathcal D\to H_{a,s}$ is continuous,
equivariant, and has image in $J$.  Assume (MA-PI.152)
and finite anticipation: there is a fixed $A_0$ such
that
\[
 \operatorname{supp}Tf\subset
       [\inf\operatorname{supp}f-A_0,\infty)
                         \quad(f\in\mathcal D).
 \tag{MA-PI.164}
\]
Then $T=0$.  The same holds on $\mathcal D_0$.
Here the proof uses Conrey's established positive
proportion of simple critical-line zeros, precisely
as recorded in MA-MO-023.  This is an explicit partial
zero-location input; this variant is not presented
as an argument without zero-location input.

\begin{proof}
Steps 1 and 2 of MA-PI-005 require only the stated
Hilbert continuity and remain valid.  By finite
anticipation and weighted Cauchy--Schwarz,
$\mathcal M(Tf)$ is holomorphic on $\Re z<a$.
The convolution identity and local quotient patching
of Step 3 give a holomorphic multiplier $m$ on that
half-plane.  Its values on $\Re z=-a$ agree with the
actual Fourier multiplier in (MA-PI.156).
Both factors in (MA-PI.160) are therefore holomorphic
on
\[
                         -3a<\Re z<a,
 \tag{MA-PI.165}
\]
and the identity holds throughout that strip.
Every critical-line zeta value functional is
continuous on $H_{a,s}$ and annihilates $J$, as in
(MA-MO.501).  Choosing a compact input nonzero at
each Mellin point proves $m(i\gamma)=0$ at every
simple critical-line zero.  The resulting distinct
zeros of $E$ at $\gamma-ia$ number at least
$cT\log T$, by the exact Conrey input of MO023.
The same Jensen contradiction proves $T=0$.
Precomposition with $D-1/2$ proves the moment-restricted
version, retaining finite anticipation.
\end{proof}

For the actual global PW ranges,
$\operatorname{Ran}P_Y\subset H_{c}$ with $c=-\log Y$
is MA-GP-003.  Consequently the inside-plateau identity
in (MA-PI.163), applied after translating a compact
input far left, already forces (MA-PI.164) with
$A_0=\log2$.  Thus no separate support hypothesis
needs to be guessed when applying this Hilbert
variant to those exact plateau identities.

\paragraph{Inputs, current gate, and Grothendieck placement.}
MA-PI-005 uses the original smooth theta divisibility
identity of MT-M18/MT-M24 and the classical full zero
count just cited, also recorded in GRH\_NIGHT NT-R08c;
it makes no narrow-strip zero-location assumption and
does not use A or M.  Its Hilbert-only variant uses
the partial critical-line theorem explicitly, with
the inherited source-statement reading scope of MO023,
not a fresh audit of Conrey's whole mollifier proof.
Every Mellin normalization and projection remains in
the displayed metric; no transfer to the draft's
different printed $W^{s,2}$ inner product is inferred.

The original unchanged geometric identity still has
the independent FT002 arithmetic deficiency gate.
The new result rules out an entire class of exact
equivariant physical-cutoff transports; it does not
rule out a changed Gram operator, a different cutoff
family, or a scalar trace comparison.  In the cote 1
framework of BR-TR-01/02, scans079--080 and111--112,
it is an obstruction before any nuclear trace pairing
or Fredholm determinant is assigned.  The full theta
divisor is the actual finite-length local-module
data of MT-M18.  No faithful positivity or motivic
weight assertion is introduced by orthogonalizing
the cutoff ranges.

Greni\'e, L., \& Molteni, G. (2019). An explicit
Chebotarev density theorem under GRH.
\emph{Journal of Number Theory, 200}, 441--485.
\url{https://doi.org/10.1016/j.jnt.2018.12.005}.
Theorem A.1 in its Appendix A is explicitly unconditional;
the conditional theorem in the paper's title is not used.
Its statement, multiplicity convention and displayed main
term were freshly checked in the author PDF; its full
appendix proof was not freshly audited here.
The separate Hilbert variant uses Conrey (1989), Theorem 1,
p.~4 and counts (8), (11), (12), p.~3, with the source
scope and full bibliographic entry recorded in MA-MO-023.
